\documentclass[../../main.tex]{subfiles}% 注意这里的文件路径不能用 ./main.tex ,否则用latexmk编译子文件会报错
\graphicspath{{\subfix{./image/}}{\subfix{../../image/}}} % 指定图片引用路径,后续可以直接使用图片文件名
% 如果图片存放在上级目录,则使用 ../../image/ 作为备用路径

% 例如:
% \begin{figure}[H]
% \centering
% \includegraphics[width=0.3\textwidth]{图.png}
% \caption{}
% \label{figure:图}
% \end{figure}
% 注意:上述\label{}一定要放在\caption{}之后,否则引用图片序号会只会显示??.

\begin{document}

\section{混合问题}

在这一节中我们将考虑一个更实际的模型。在现实生活中我们研究的弹性体都是有界的。由于这些物体所在的区域具有边界，边界点的位移或受力情况会通过弹性应力对弹性体产生影响，因此我们对波动方程既需要提初始条件，又需要提边值条件。这样的定解问题通常称为混合问题。

我们重点考虑一维混合问题，其中给定的边值条件是第一类的。实际模型就是一根有限长的振动的均匀细弦。我们已知弦上各点的初位移与初速度和受力情况，并且知道两端点的位移变化，要求在振动过程中弦上各点的位移。我们将利用分离变量法来求解这样的混合问题。对于其他边值条件的情形我们也可以利用同样的方法求解。事实上分离变量法也适用于求解一些多维混合问题，但求解过程要复杂一些。

\subsection{分离变量法}

现在我们考虑这样的混合问题
\begin{align}
\begin{cases}
u_{tt}-a^2u_{xx}=f, & (x,t)\in(0,l)\times(0,+\infty), \\
u(x,0)=\varphi(x), & x\in[0,l], \\
u_t(x,0)=\psi(x), & x\in[0,l], \\
u(0,t)=g_1(t),\quad u(l,t)=g_2(t), & t\in[0,+\infty).
\end{cases} \label{eq:wave-mixed}
\end{align}
我们将利用分离变量法来求解上述混合问题。为了更好地理解分离变量法，我们从最简单的情形开始。

\subsubsection{$f\equiv 0,\ g_1(t)=g_2(t)\equiv 0$ 的情形}

现在我们考虑混合问题
\begin{align}
\begin{cases}
u_{tt}-a^2u_{xx}=0, & (x,t)\in(0,l)\times(0,+\infty), \\
u(x,0)=\varphi(x), & x\in[0,l], \\
u_t(x,0)=\psi(x), & x\in[0,l], \\
u(0,t)=u(l,t)=0, & t\in[0,+\infty).
\end{cases} \label{eq:wave-mixed-hom}
\end{align}
我们利用分离变量法来求解此问题。

具体来说，考虑分离变量形式的非零解
\begin{align*}
u(x,t)=X(x)T(t).
\end{align*}
将它代入混合问题 \eqref{eq:wave-mixed-hom} 的方程得到
\begin{align*}
T''(t)X(x)-a^2X''(x)T(t)=0,\quad (x,t)\in(0,l)\times(0,+\infty),
\end{align*}
即
\begin{align*}
\frac{T''(t)}{a^2T(t)}=\frac{X''(x)}{X(x)}.
\end{align*}
在上式中，左端是 $t$ 的函数，右端是 $x$ 的函数，因而只能是常数，记为 $-\lambda$，从而
\begin{align*}
T''(t)+a^2\lambda T(t)=0,\quad t\in[0,+\infty),\\
X''(x)+\lambda X(x)=0,\quad x\in(0,l).
\end{align*}
将 $u(x,t)=X(x)T(t)$ 代入混合问题 \eqref{eq:wave-mixed-hom} 的边值条件得到
\begin{align*}
X(0)T(t)=X(l)T(t)=0,\quad t\in(0,+\infty).
\end{align*}
由于 $u(x,t)\neq 0$，因此 $T(t)\neq 0$，从而
\begin{align*}
X(0)=X(l)=0.
\end{align*}
于是我们得到特征值问题
\begin{align}
\begin{cases}
X''(x)+\lambda X(x)=0, & x\in(0,l), \\
X(0)=X(l)=0.
\end{cases} \label{eq:wave-eigen}
\end{align}
根据定理 \eqref{eq:sturm-liouville} 的结论，易见特征值问题 \eqref{eq:wave-eigen} 的解必满足
\begin{align*}
X(x)=C_1\cos\sqrt{\lambda}x+C_2\sin\sqrt{\lambda}x,
\end{align*}
其中 $C_1,C_2$ 为待定常数，由边值条件得到
\begin{align*}
C_1=0,\qquad C_2\sin\sqrt{\lambda}l=0.
\end{align*}
由于 $C_2\neq 0$，否则 $X(x)\equiv 0$，从而
\begin{align*}
\sin\sqrt{\lambda}l=0,
\end{align*}
于是所有特征值是
\begin{align*}
\lambda_n=\left(\frac{n\pi}{l}\right)^2,\quad n=1,2,\cdots,
\end{align*}
与特征值 $\lambda_n$ 对应的特征函数是
\begin{align*}
X_n(x)=\sin\frac{n\pi}{l}x,\quad n=1,2,\cdots.
\end{align*}
而与 $X_n(x)$ 对应的 $T_n(t)$ $(n=1,2,\cdots)$ 为
\begin{align*}
T_n(t)=A_n\cos\frac{na\pi}{l}t+B_n\sin\frac{na\pi}{l}t,\quad n=1,2,\cdots,
\end{align*}
其中 $A_n,B_n$ $(n=1,2,\cdots)$ 为待定常数。分离变量形式的解为
\begin{align*}
u_n(x,t)=T_n(t)X_n(x)=\left(A_n\cos\frac{na\pi}{l}t+B_n\sin\frac{na\pi}{l}t\right)\sin\frac{n\pi}{l}x,\quad n=1,2,\cdots.
\end{align*}
从形式上看，每一个 $u_n(x,t)$ $(n=1,2,\cdots)$ 都满足方程和边值条件，但一般来讲它们都不满足初始条件。为求一个既满足方程和边值条件又满足初始条件的解，我们将 $u_n(x,t)$ 叠加。形式上，
\begin{align}
u(x,t)=\sum_{n=1}^{\infty}u_n(x,t)
=\sum_{n=1}^{\infty}\left(A_n\cos\frac{na\pi}{l}t+B_n\sin\frac{na\pi}{l}t\right)\sin\frac{n\pi}{l}x \label{eq:wave-sep-sol}
\end{align}
满足方程。为使 $u(x,t)$ 满足初始条件我们需要
\begin{align*}
u(x,0)=\varphi(x)=\sum_{n=1}^{\infty}A_n\sin\frac{n\pi}{l}x,\\
u_t(x,0)=\psi(x)=\sum_{n=1}^{\infty}\frac{na\pi}{l}B_n\sin\frac{n\pi}{l}x.
\end{align*}
由于特征函数系 $\left\{\sin\frac{n\pi}{l}x\right\}\ (n=1,2,\cdots)$ 的完备性，因此 $\varphi(x)$，$\psi(x)$ 可以展开成
\begin{align*}
\varphi(x)=\sum_{n=1}^{\infty}\varphi_n\sin\frac{n\pi}{l}x,\qquad
\psi(x)=\sum_{n=1}^{\infty}\psi_n\sin\frac{n\pi}{l}x,
\end{align*}
其中
\begin{align*}
\varphi_n=\frac{2}{l}\int_0^l\varphi(x)\sin\frac{n\pi}{l}x\,\mathrm{d}x,\quad n=1,2,\cdots,\\
\psi_n=\frac{2}{l}\int_0^l\psi(x)\sin\frac{n\pi}{l}x\,\mathrm{d}x,\quad n=1,2,\cdots.
\end{align*}
比较上述表达式我们得到
\begin{align}
A_n=\varphi_n=\frac{2}{l}\int_0^l\varphi(x)\sin\frac{n\pi}{l}x\,\mathrm{d}x,\quad n=1,2,\cdots, \label{eq:wave-A}
\end{align}
\begin{align}
B_n=\frac{l}{na\pi}\psi_n=\frac{2}{na\pi}\int_0^l\psi(x)\sin\frac{n\pi}{l}x\,\mathrm{d}x,\quad n=1,2,\cdots. \label{eq:wave-B}
\end{align}
这样我们就求得混合问题 \eqref{eq:wave-mixed-hom} 的形式解 \eqref{eq:wave-sep-sol}，其中它的系数由式 \eqref{eq:wave-A}, \eqref{eq:wave-B} 决定。为了使得所得到的形式解确实是混合问题 \eqref{eq:wave-mixed-hom} 的古典解，我们需要对初值提适当的光滑性要求，同时需要在定解区域的角点 $(0,0)$，$(l,0)$ 提适当的相容性条件。这里相容性条件的导出与上一小节半无界问题的相容性条件完全一致，参见习题。

\begin{theorem}
若 $\varphi(x)\in C^3[0,l]$，$\psi(x)\in C^2[0,l]$ 及 $\varphi(x),\psi(x)$ 在定解区域 $Q=(0,l)\times(0,+\infty)$ 的角点 $(0,0)$，$(l,0)$ 满足相容性条件
\begin{align*}
\varphi(0)=\varphi(l)=\varphi''(0)=\varphi''(l)=\psi(0)=\psi(l)=0,
\end{align*}
则由式 \eqref{eq:wave-sep-sol} 给出的函数 $u(x,t)\in C^2(\overline{Q})$，且满足混合问题 \eqref{eq:wave-mixed-hom}。
\end{theorem}
\begin{proof}
由定理的假设和分部积分公式，我们得到
\begin{align*}
A_n=\frac{2}{l}\int_0^l\varphi(x)\sin\frac{n\pi}{l}x\,\mathrm{d}x=-\frac{l^3}{\pi^3n^3}a_n,\quad n=1,2,\cdots,
\end{align*}
其中
\begin{align*}
a_n=\frac{2}{l}\int_0^l\varphi'''(x)\sin\frac{n\pi}{l}x\,\mathrm{d}x,\quad n=1,2,\cdots.
\end{align*}
同理
\begin{align*}
B_n=\frac{2}{na\pi}\int_0^l\psi(x)\sin\frac{n\pi}{l}x\,\mathrm{d}x=-\frac{l^3}{a\pi^3n^3}b_n,\quad n=1,2,\cdots,
\end{align*}
其中
\begin{align*}
b_n=\frac{2}{l}\int_0^l\psi''(x)\sin\frac{n\pi}{l}x\,\mathrm{d}x,\quad n=1,2,\cdots.
\end{align*}
于是在 $\overline{Q}=[0,l]\times[0,+\infty)$ 上我们有估计
\begin{align*}
|u_n(x,t)|&\leqslant |A_n|+|B_n|\leqslant \frac{C_1}{n^3},\quad n=1,2,\cdots,\\
|Du_n(x,t)|&\leqslant \frac{C_2}{n^2},\quad n=1,2,\cdots,\\
|D^2u_n(x,t)|&\leqslant a_n^2+b_n^2+\frac{C_3}{n^2},\quad n=1,2,\cdots,
\end{align*}
这里 $Du_n(x,t)$ 表示 $u_n(x,t)$ 的所有一阶偏导数，$|Du_n(x,t)|$ 表示 $u_n(x,t)$ 的所有一阶偏导数的平方和的平方根，$D^2u_n(x,t)$ 表示 $u_n(x,t)$ 的所有二阶偏导数，$|D^2u_n(x,t)|$ 表示 $u_n(x,t)$ 的所有二阶偏导数的平方和的平方根，$C_1,C_2,C_3$ 是只依赖于 $l,a$ 和定积分
\begin{align*}
\int_0^l|\varphi''''(x)|\,\mathrm{d}x,\qquad \int_0^l|\psi'''(x)|\,\mathrm{d}x
\end{align*}
的常数。

从 Fourier 级数的 Bessel 不等式得到
\begin{align*}
\sum_{n=1}^{\infty}a_n^2&\leqslant \frac{2}{l}\int_0^l|\varphi'''(x)|^2\,\mathrm{d}x,\\
\sum_{n=1}^{\infty}b_n^2&\leqslant \frac{2}{l}\int_0^l|\psi''(x)|^2\,\mathrm{d}x.
\end{align*}
因此我们证明了级数
\begin{align*}
\sum_{n=1}^{\infty}u_n(x,t),\qquad \sum_{n=1}^{\infty}Du_n(x,t),\qquad \sum_{n=1}^{\infty}D^2u_n(x,t)
\end{align*}
在 $\overline{Q}$ 上均一致收敛，从而得到由表达式 \eqref{eq:wave-sep-sol} 给出的函数 $u\in C^2(\overline{Q})$。而且我们可以将 $u(x,t)$ 的级数逐项对变量 $x,t$ 微分两次，并容易证明 $u(x,t)$ 满足混合问题 \eqref{eq:wave-mixed-hom} 的方程和初边值条件。
\end{proof}

\subsubsection{$f(x,t)\neq 0,\ g_1(t)=g_2(t)=0$ 的情形}

现在我们考虑混合问题
\begin{align}
\begin{cases}
u_{tt}-a^2u_{xx}=f, & (x,t)\in(0,l)\times(0,+\infty), \\
u(x,0)=\varphi(x), & x\in[0,l], \\
u_t(x,0)=\psi(x), & x\in[0,l], \\
u(0,t)=u(l,t)=0, & t\in[0,+\infty).
\end{cases} \label{eq:wave-mixed-nonhom}
\end{align}
我们仍然可以利用分离变量法来求解。

具体来说，把混合问题 \eqref{eq:wave-mixed-nonhom} 的解 $u(x,t)$，非齐次项 $f(x,t)$ 和初值 $\varphi(x)$，$\psi(x)$ 都按特征函数系 $\left\{\sin\frac{n\pi}{l}x\right\}\ (n=1,2,\cdots)$ 展开：
\begin{align}
u(x,t)=\sum_{n=1}^{\infty}T_n(t)\sin\frac{n\pi}{l}x, \label{eq:wave-u-series}
\end{align}
\begin{align*}
f(x,t)=\sum_{n=1}^{\infty}f_n(t)\sin\frac{n\pi}{l}x,\\
\varphi(x)=\sum_{n=1}^{\infty}\varphi_n\sin\frac{n\pi}{l}x,\\
\psi(x)=\sum_{n=1}^{\infty}\psi_n\sin\frac{n\pi}{l}x.
\end{align*}
由特征函数系 $\left\{\sin\frac{n\pi}{l}x\right\}$ 的正交性和完备性得到
\begin{align*}
f_n(t)=\frac{2}{l}\int_0^l f(\xi,t)\sin\frac{n\pi}{l}\xi\,\mathrm{d}\xi,\quad n=1,2,\cdots,\\
\varphi_n=\frac{2}{l}\int_0^l\varphi(\xi)\sin\frac{n\pi}{l}\xi\,\mathrm{d}\xi,\quad n=1,2,\cdots,\\
\psi_n=\frac{2}{l}\int_0^l\psi(\xi)\sin\frac{n\pi}{l}\xi\,\mathrm{d}\xi,\quad n=1,2,\cdots.
\end{align*}
为求出未知函数 $T_n(t)$，把上述表达式代入混合问题 \eqref{eq:wave-mixed-nonhom} 的方程，由特征函数系的完备性，从而得到 $T_n(t)$ 满足以下常微分方程和初值条件
\begin{align}
\begin{cases}
T_n''(t)+\left(\frac{n\pi a}{l}\right)^2T_n(t)=f_n(t), & t\in(0,+\infty), \\
T_n(0)=\varphi_n,\quad T_n'(0)=\psi_n.
\end{cases} \label{eq:wave-ode}
\end{align}
求解常微分方程初值问题 \eqref{eq:wave-ode} 得到，对于 $n=1,2,\cdots$，
\begin{align}
T_n(t)=\varphi_n\cos\frac{n\pi a}{l}t+\frac{l}{n\pi a}\psi_n\sin\frac{n\pi a}{l}t
+\frac{l}{n\pi a}\int_0^t f_n(\tau)\sin\frac{n\pi a}{l}(t-\tau)\,\mathrm{d}\tau. \label{eq:wave-Tn}
\end{align}
将它代入级数 \eqref{eq:wave-u-series}，这样就求得混合问题 \eqref{eq:wave-mixed-nonhom} 的解的表达式。如果对初值 $\varphi,\psi$ 和非齐次项 $f$ 提适当的光滑性要求，并在角点 $(0,0)$，$(l,0)$ 提相容性条件
\begin{align*}
\varphi(0)=\varphi(l)=\psi(0)=\psi(l)=0,\\
a^2\varphi''(0)+f(0,0)=0,\qquad a^2\varphi''(l)+f(l,0)=0,
\end{align*}
同样可以证明上面的表达式给出的形式解确实是混合问题 \eqref{eq:wave-mixed-nonhom} 属于 $C^2(\overline{Q})$ 的解。

\subsubsection{$f(x,t)\neq 0,\ g_1(t)\neq 0,\ g_2(t)\neq 0$ 的情形}

现在我们考虑具有非齐次边值的混合问题
\begin{align}
\begin{cases}
u_{tt}-a^2u_{xx}=f, & (x,t)\in(0,l)\times(0,+\infty), \\
u(x,0)=\varphi(x), & x\in[0,l], \\
u_t(x,0)=\psi(x), & x\in[0,l], \\
u(0,t)=g_1(t),\quad u(l,t)=g_2(t), & t\in[0,+\infty).
\end{cases} \label{eq:wave-mixed-nonhom-bc}
\end{align}
为了利用分离变量法我们作函数变换
\begin{align*}
v=u-\left[\frac{l-x}{l}g_1(t)+\frac{x}{l}g_2(t)\right],
\end{align*}
得到一个关于 $v$ 的具有齐次边值的混合问题。对初值 $\varphi,\psi$，边值 $g_1,g_2$，和非齐次项 $f$ 提适当的光滑性要求，在角点 $(0,0)$，$(l,0)$ 提相容性条件
\begin{gather*}
\varphi(0)=g_1(0),\quad \varphi(l)=g_2(0),\\
\psi(0)=g_1'(0),\quad \psi(l)=g_2'(0),\\
g_1''(0)-a^2\varphi''(0)=f(0,0),\quad g_2''(0)-a^2\varphi''(l)=f(l,0),
\end{gather*}
我们仍然可以证明所得的解是混合问题的古典解。

\begin{remark}
对于混合问题 \eqref{eq:wave-mixed-hom}，我们同样可以采用对称开拓法求解。具体来说，对初值 $\varphi,\psi$ 进行奇且周期为 $2l$ 的对称开拓，定义
\begin{align*}
\Phi(x)=\varphi(x),\quad x\in[0,l],\qquad \Phi(x)=-\Phi(-x),\quad \Phi(x)=\Phi(x+2l),\quad x\in\mathbb{R},\\
\Psi(x)=\psi(x),\quad x\in[0,l],\qquad \Psi(x)=-\Psi(-x),\quad \Psi(x)=\Psi(x+2l),\quad x\in\mathbb{R}.
\end{align*}
然后求解初值问题
\begin{align*}
\begin{cases}
U_{tt}-a^2U_{xx}=0, & (x,t)\in\mathbb{R}\times\mathbb{R}_+, \\
U(x,0)=\Phi(x), & x\in\mathbb{R}, \\
U_t(x,0)=\Psi(x), & x\in\mathbb{R}.
\end{cases}
\end{align*}
由 D'Alembert 公式，得到解
\begin{align*}
U(x,t)=\frac{1}{2}[\Phi(x+at)+\Phi(x-at)]+\frac{1}{2a}\int_{x-at}^{x+at}\Psi(\xi)\,\mathrm{d}\xi.
\end{align*}
容易证明 $U(x,t)$ 是关于 $x$ 周期为 $2l$ 的奇函数，因此在 $\overline{Q}$ 上，令
\begin{align*}
u(x,t)=U(x,t),
\end{align*}
则它必然是混合问题 \eqref{eq:wave-mixed-hom} 的解。由上述解的表达式可以看出，定理中关于初值 $\varphi$ 和 $\psi$ 的光滑性要求可以降低为 $\varphi(x)\in C^2[0,l]$，$\psi(x)\in C^1[0,l]$。
\end{remark}

在历史上，解的表达式 \eqref{eq:wave-sep-sol} 首先是从物理直观得到的。从解的表达式 \eqref{eq:wave-sep-sol} 我们知道混合问题 \eqref{eq:wave-mixed-hom} 的解可以表示成三角级数的和，但是人们同时又利用对称开拓法把解表示成 D'Alembert 公式。这两个解的表达式差别如此大，这是什么原因呢？实际上，如果对 Fourier 级数有深入的了解，我们将知道上述两个解的表达式表示同样的函数。此结论可以利用 Fourier 级数直接证明，也可以根据后面的能量不等式来证明。这是因为能量不等式蕴涵着混合问题的古典解的唯一性，从而证明两个不同的表达式表示同样的函数。

\subsection{驻波法与共振}

分离变量法具有明显的物理意义。实际上，我们把级数 \eqref{eq:wave-sep-sol} 中的每一项 $u_n(x,t)$ 改写成
\begin{align*}
u_n(x,t)=N_n\sin\frac{n\pi}{l}x\sin\left(\frac{na\pi}{l}t+\alpha_n\right),
\end{align*}
其中 $N_n=\sqrt{A_n^2+B_n^2}$，$\alpha_n=\arctan\frac{A_n}{B_n}$。$u_n(x,t)$ 称为\textbf{弦的振动元素}。描述了弦上的点 $x$ 所做的振幅为 $a_n=N_n\sin\frac{n\pi}{l}x$，频率为 $\omega_n=\frac{na\pi}{l}$，初相位为 $\alpha_n$ 的简谐振动。特别地，如果初相位 $\alpha_n=0$，当 $x=0,\frac{l}{n},\cdots,\frac{n-1}{n}l,l$ 时，振动元素 $u_n(x,t)$ 的振幅 $a_n=0$；当 $x=\frac{l}{2n},\frac{3l}{2n},\cdots,\frac{2n-1}{2n}l$ 时，$u_n(x,t)$ 的振幅 $a_n=\pm N_n$。弦的这种形式的运动称为\textbf{驻波}。因此在物理学上分离变量法也称为\textbf{驻波法}。由分离变量法我们知道，弦的振动在不受外力作用下，可以分解为一系列具有特定频率的驻波的叠加。

下面我们利用一维波动方程的混合问题的结论来解释弦乐器（小提琴、琵琶、吉他、二胡等）的演奏原理。我们用 $u(x,t)$ 表示弦乐器的弦的振动位移。由于弦的振动而引起其周围空气的振动，进一步引起音箱的共鸣，从而产生出悦耳的音乐。因此我们可以用由分离变量法得到的解的表达式 $u(x,t)=\sum_{n=1}^{\infty}u_n(x,t)$ 来表示我们听到的音乐。此时弦的基音是由最低频率 $\omega_1=\frac{a\pi}{l}=\frac{\pi}{l}\sqrt{\frac{T_0}{\rho_0}}$（$T_0$ 为弦的张力，$\rho_0$ 为密度）所对应的单音 $u_1(x,t)$ 确定的。振动元素 $u_1(x,t)$ 的最大振幅 $N_1=\sqrt{A_1^2+B_1^2}$ 通常要比振动元素 $u_n(x,t)$ 的最大振幅 $N_n=\sqrt{A_n^2+B_n^2}$ $(n\geqslant 2)$ 大得多，因此它决定了音乐的基调。而弦的其余的单音 $u_n(x,t)$ $(n\geqslant 2)$ 称为泛音，它们构成了声音的音色。我们知道不同的弦乐器演奏同一首曲子时，虽然音调是相同的，但是音乐声却是不同的。这是因为它们虽然具有相同的基音频率，却也具有不同的泛音频率。这样就产生了音质的差异。在用弦乐器演奏曲子的时候，由于演奏者用手指按住弦线的不同部位，这样受振动的弦的长度顿时变小了，基音频率 $\omega_1=\frac{\pi}{l}\sqrt{\frac{T_0}{\rho_0}}$ 就增大，音调也随之升高。特别地，当手指按住弦的中点时，基音和泛音的频率就比原来的频率增加一倍，弦就发出比原来高八度的音调。另外，演奏者也经常利用拧紧和拧松弦线的方法来调整弦的音调。这也可以从弦的基音频率的表达式来理解。当拧紧弦线时，弦的张力增大，弦的音调就变高；当拧松弦线时，弦的张力减小，弦的音调就变低。此外，我们知道每件弦乐器上都有好几根长度相同而粗细不同的弦，而且每根弦演奏的音调也不同。从弦的基音频率的表达式来看，由于粗弦的密度大，因而基音频率就小，音调就低；由于细弦的密度小，因而基音频率就大，音调就高。从上面的分析我们就对弦乐器的演奏原理得到初步的理解。

由用分离变量法得到的混合问题 \eqref{eq:wave-mixed-nonhom} 的解的表达式 \eqref{eq:wave-u-series} 和 \eqref{eq:wave-Tn}，我们还可以解释物理学中在强迫振动下所产生的\textbf{共振现象}。为此考虑混合问题
\begin{align*}
\begin{cases}
u_{tt}-a^2u_{xx}=A(x)\sin\omega t, & (x,t)\in(0,l)\times(0,+\infty), \\
u(x,0)=0, & x\in[0,l], \\
u_t(x,0)=0, & x\in[0,l], \\
u(0,t)=u(l,t)=0, & t\in[0,+\infty),
\end{cases}
\end{align*}
其中 $\omega$ 是一个正常数，函数 $A(x)\in C[0,l]$ 满足 $A(0)=A(l)=0$。
由解的表达式 \eqref{eq:wave-u-series} 知
\begin{align*}
u(x,t)&=\sum_{n=1}^{\infty}\sin\frac{n\pi}{l}x\left[\frac{l}{na\pi}\int_0^t f_n(\tau)\sin\frac{na\pi}{l}(t-\tau)\,\mathrm{d}\tau\right] \\
&=\sum_{n=1}^{\infty}\frac{a_n}{\omega_n}\sin\frac{n\pi}{l}x\int_0^t\sin\omega\tau\sin\omega_n(t-\tau)\,\mathrm{d}\tau,
\end{align*}
其中
\begin{align*}
\omega_n=\frac{n\pi a}{l},\qquad a_n=\frac{2}{l}\int_0^l A(x)\sin\frac{n\pi}{l}x\,\mathrm{d}x.
\end{align*}
当 $\omega\neq\omega_n$ $(n=1,2,\cdots)$ 时，
\begin{align*}
u(x,t)=\sum_{n=1}^{\infty}\frac{a_n}{\omega_n(\omega^2-\omega_n^2)}(\omega\sin\omega_nt-\omega_n\sin\omega t)\sin\frac{n\pi}{l}x.
\end{align*}
显然上等式右端的级数是一致有界的。

当存在某个 $k\in\mathbb{Z}$ 使得 $\omega=\omega_k$ 时，则
\begin{align*}
u(x,t)&=\sum_{n\neq k}^{\infty}\frac{a_n}{\omega_n(\omega^2-\omega_n^2)}(\omega_k\sin\omega_nt-\omega_n\sin\omega_kt)\sin\frac{n\pi}{l}x \\
&\quad+\frac{a_k}{2\omega_k^2}(\sin\omega_kt-t\omega_k\cos\omega_kt)\sin\frac{k\pi}{l}x.
\end{align*}
此时上等式右端的级数是一致有界的，而当选取 $A(x)$ 使得 $a_k\neq 0$ 时，对应于固有频率 $\omega_k$ 的第 $k$ 个振动元素 $u_k(x,t)$ 的振幅可以随时间一起无限增大，这就是物理上的共振现象。这表示一根两端固定的弦在一个周期外力的作用下做强迫振动，如果周期外力的频率与弦的某个固有频率相等，那么弦将产生共振，弦的一些点的振幅将随着时间的增大而趋于无穷，这必然在某一时刻导致弦的断裂。

实际上，此例还蕴涵着在位势方程和热方程中占有重要地位的极大模估计在波动方程中不可能成立。这也反映了波动方程的解的性质与位势方程和热方程的解的性质有很大的差别。

\subsection{能量不等式}

\begin{theorem}
设 $\Omega$ 是 $\mathbb{R}^n$ 上的一个有界开集，并令 $\Omega_T=\Omega\times(0,T)$。如果 $u\in C^1(\overline{\Omega}_T)\cap C^2(\Omega_T)$ 是混合问题
\begin{align}
\begin{cases}
u_{tt}-a^2\Delta u=f, & (x,t)\in\Omega\times(0,T), \\
u(x,0)=\varphi(x), & x\in\overline{\Omega}, \\
u_t(x,0)=\psi(x), & x\in\overline{\Omega}, \\
u(x,t)=0, & t\in[0,T],\ x\in\partial\Omega,
\end{cases} \label{eq:wave-mixed-energy-gen}
\end{align}
的解，则存在只依赖于 $T$ 的常数 $M$，使得对于任意 $\tau\in[0,T]$，有
\begin{align}
\int_\Omega [u_t^2(x,\tau)+a^2|Du(x,\tau)|^2]\,\mathrm{d}x
\leqslant M\left\{\int_\Omega[\psi^2(x)+a^2|D\varphi(x)|^2]\,\mathrm{d}x+\iint_{\Omega_\tau}f^2(x,t)\,\mathrm{d}x\,\mathrm{d}t\right\}, \label{eq:energy-est-wave}
\end{align}
其中 $\Omega_\tau=\Omega\times(0,\tau)$。当 $f\equiv 0$ 时，$M=1$ 且上式中的不等号由等号代替。
\end{theorem}

为简单起见，我们只对二维情形给出证明。此时定理具有如下形式。

\begin{theorem}
设 $\Omega$ 是 $\mathbb{R}^2$ 上的一个有界开集，$\Omega_T=\Omega\times(0,T)$。如果 $u\in C^1(\overline{\Omega}_T)\cap C^2(\Omega_T)$ 是混合问题 \eqref{eq:wave-mixed-energy-gen} 的解，则存在只依赖于 $T$ 的常数 $M$，使得对于 $\tau\in[0,T]$，有
\begin{align}
&\iint_\Omega [u_t^2(x,y,\tau)+a^2(u_x^2(x,y,\tau)+u_y^2(x,y,\tau))]\,\mathrm{d}x\mathrm{d}y
\nonumber \\
&\leqslant M\Bigg\{\iint_\Omega[\psi^2(x,y)+a^2(\varphi_x^2(x,y)+\varphi_y^2(x,y))]\,\mathrm{d}x\mathrm{d}y 
+\iiint_{\Omega_\tau}f^2(x,y,t)\,\mathrm{d}x\mathrm{d}y\mathrm{d}t\Bigg\}, \label{eq:energy-est-2d}
\end{align}
其中 $\Omega_\tau=\Omega\times(0,\tau)$（见\reffig{figure:omega_tau}）。当 $f\equiv 0$ 时，$M=1$ 且上式中的不等号由等号代替。
\end{theorem}

\begin{figure}[H]
\centering
\includegraphics[width=0.5\textwidth]{PDE图6.png}
\caption{区域 $\Omega_\tau$ 示意}
\label{figure:omega_tau}
\end{figure}

\begin{proof}
在波动方程两端同乘 $u_t$，并在区域 $\Omega_\tau$ 上积分，得到
\begin{align*}
\iiint_{\Omega_\tau}u_t[u_{tt}-a^2(u_{xx}+u_{yy})]\,\mathrm{d}x\mathrm{d}y\mathrm{d}t=\iiint_{\Omega_\tau}u_tf\,\mathrm{d}x\mathrm{d}y\mathrm{d}t.
\end{align*}
注意到
\begin{align*}
u_t[u_{tt}-a^2(u_{xx}+u_{yy})]
=\frac{1}{2}[u_t^2+a^2(u_x^2+u_y^2)]_t-a^2[(u_tu_x)_x+(u_tu_y)_y],
\end{align*}
将左端记为 $I$，应用 Gauss-Green 公式得
\begin{align*}
I=\iint_{\partial\Omega_\tau}\left\{\frac{1}{2}[u_t^2+a^2(u_x^2+u_y^2)]\cos(n,t)-a^2[u_tu_x\cos(n,x)+u_tu_y\cos(n,y)]\right\}\mathrm{d}S.
\end{align*}
在柱体 $\Omega\times[0,\tau]$ 的侧面 $\partial\Omega\times[0,\tau]$ 上，$\cos(n,t)=0$，且由边值条件 $u|_{\partial\Omega}=0$ 得 $u_t|_{\partial\Omega}=0$，故侧面上的积分为零。因此
\begin{align*}
I=\frac{1}{2}\iint_\Omega [u_t^2(x,y,\tau)+a^2(u_x^2(x,y,\tau)+u_y^2(x,y,\tau))]\,\mathrm{d}x\mathrm{d}y
-\frac{1}{2}\iint_\Omega[\psi^2+a^2(\varphi_x^2+\varphi_y^2)]\,\mathrm{d}x\mathrm{d}y.
\end{align*}
于是得到
\begin{align}
\frac{1}{2}\iint_\Omega [u_t^2+a^2(u_x^2+u_y^2)]\,\mathrm{d}x\mathrm{d}y-\frac{1}{2}\iint_\Omega[\psi^2+a^2(\varphi_x^2+\varphi_y^2)]\,\mathrm{d}x\mathrm{d}y
=\iiint_{\Omega_\tau}u_tf\,\mathrm{d}x\mathrm{d}y\mathrm{d}t. \label{eq:energy-identity}
\end{align}
当 $f\equiv 0$ 时，上式右端为零，即得等号且 $M=1$。

对于一般情形，由不等式 $2ab\leqslant \varepsilon a^2+\frac{1}{\varepsilon}b^2$（$\varepsilon>0$）得
\begin{align*}
\iiint_{\Omega_\tau}u_tf\,\mathrm{d}x\mathrm{d}y\mathrm{d}t
\leqslant \frac{\varepsilon}{2}\iiint_{\Omega_\tau}u_t^2\,\mathrm{d}x\mathrm{d}y\mathrm{d}t+\frac{1}{2\varepsilon}\iiint_{\Omega_\tau}f^2\,\mathrm{d}x\mathrm{d}y\mathrm{d}t.
\end{align*}
代入 \eqref{eq:energy-identity} 得
\begin{align*}
&\iint_\Omega [u_t^2(x,y,\tau)+a^2(u_x^2(x,y,\tau)+u_y^2(x,y,\tau))]\,\mathrm{d}x\mathrm{d}y \\
&\leqslant \iint_\Omega[\psi^2+a^2(\varphi_x^2+\varphi_y^2)]\,\mathrm{d}x\mathrm{d}y
+\varepsilon\iiint_{\Omega_\tau}u_t^2\,\mathrm{d}x\mathrm{d}y\mathrm{d}t+\frac{1}{\varepsilon}\iiint_{\Omega_\tau}f^2\,\mathrm{d}x\mathrm{d}y\mathrm{d}t.
\end{align*}
记
\begin{align*}
G(\tau)=\iiint_{\Omega_\tau}[u_t^2+a^2(u_x^2+u_y^2)]\,\mathrm{d}x\mathrm{d}y\mathrm{d}t,
\end{align*}
则
\begin{align*}
\frac{\mathrm{d}G(\tau)}{\mathrm{d}\tau}=\iint_\Omega [u_t^2(x,y,\tau)+a^2(u_x^2(x,y,\tau)+u_y^2(x,y,\tau))]\,\mathrm{d}x\mathrm{d}y.
\end{align*}
于是得到微分不等式
\begin{align*}
\frac{\mathrm{d}G(\tau)}{\mathrm{d}\tau}\leqslant \varepsilon G(\tau)+F_\varepsilon(\tau),
\end{align*}
其中
\begin{align*}
F_\varepsilon(\tau)=\iint_\Omega[\psi^2+a^2(\varphi_x^2+\varphi_y^2)]\,\mathrm{d}x\mathrm{d}y+\frac{1}{\varepsilon}\iiint_{\Omega_\tau}f^2\,\mathrm{d}x\mathrm{d}y\mathrm{d}t.
\end{align*}
$F_\varepsilon(\tau)$ 是 $\tau$ 的单调递增函数。由 Gronwall 不等式（引理 \eqref{eq:gronwall-diff}）得
\begin{align*}
G(\tau)\leqslant \int_0^\tau e^{\varepsilon(\tau-t)}F_\varepsilon(t)\,\mathrm{d}t\leqslant \frac{1}{\varepsilon}(e^{\varepsilon\tau}-1)F_\varepsilon(\tau),
\end{align*}
进而
\begin{align*}
\frac{\mathrm{d}G(\tau)}{\mathrm{d}\tau}\leqslant e^{\varepsilon\tau}F_\varepsilon(\tau)\leqslant e^{\varepsilon T}F_\varepsilon(\tau).
\end{align*}
取 $\varepsilon=\frac{1}{T}$，得
\begin{align*}
\iint_\Omega [u_t^2+a^2(u_x^2+u_y^2)]\,\mathrm{d}x\mathrm{d}y
\leqslant e(1+T)\left\{\iint_\Omega[\psi^2+a^2(\varphi_x^2+\varphi_y^2)]\,\mathrm{d}x\mathrm{d}y+\iiint_{\Omega_\tau}f^2\,\mathrm{d}x\mathrm{d}y\mathrm{d}t\right\}.
\end{align*}
定理获证。
\end{proof}

\begin{remark}
若不记常数因子，表达式
\begin{align*}
\iint_\Omega [u_t^2(x,y,\tau)+a^2(u_x^2(x,y,\tau)+u_y^2(x,y,\tau))]\,\mathrm{d}x\mathrm{d}y
\end{align*}
表示物体 $\Omega$ 在 $\tau$ 时刻的总能量，其中第一项为动能，第二项为弹性势能。当 $f\equiv 0$ 时，总能量守恒。
\end{remark}

\begin{remark}
对不等式 \eqref{eq:energy-est-2d} 关于 $\tau$ 从 $0$ 到 $t$ 积分，可得到另一种形式的能量模估计
\begin{align}
\iiint_{\Omega_t}[u_t^2+a^2(u_x^2+u_y^2)]\,\mathrm{d}x\mathrm{d}y\mathrm{d}t
\leqslant M\left\{\iint_\Omega[\psi^2+a^2(\varphi_x^2+\varphi_y^2)]\,\mathrm{d}x\mathrm{d}y+\iiint_{\Omega_t}f^2\,\mathrm{d}x\mathrm{d}y\mathrm{d}t\right\}, \label{eq:energy-est-int}
\end{align}
其中 $M$ 只依赖于 $T$，$\Omega_t=\Omega\times(0,t)$。同样可由此得到解 $u$ 的 $L^2$ 估计。
\end{remark}

\begin{remark}
从能量模估计 \eqref{eq:energy-est-wave} 不难导出混合问题解在函数类 $C^1(\overline{\Omega}_T)\cap C^2(\Omega_T)$ 中的惟一性以及在能量模意义下解对于初值和右端项的连续依赖性。
\end{remark}


\subsection{广义解}

在这一小节中我们考虑下列较为简单的混合问题
\begin{align}
\begin{cases}
u_{tt}-a^2u_{xx}=0, & (x,t)\in(0,l)\times(0,+\infty), \\
u(x,0)=\varphi(x), & x\in[0,l], \\
u_t(x,0)=\psi(x), & x\in[0,l], \\
u(0,t)=u(l,t)=0, & t\in[0,+\infty).
\end{cases} \label{eq:wave-gen-mixed}
\end{align}

在前面对上述混合问题的求解过程中，我们都要求所找的解在定解区域 $Q_T=(0,l)\times(0,T)$ 的内部二次连续可微且满足方程，同时在区域边界上满足初边值条件。为了保证混合问题 \eqref{eq:wave-gen-mixed} 在函数类 $C^1(\overline{Q}_T)\cap C^2(\overline{Q}_T)$ 有解，我们对于初边值提了一定的光滑性条件，并在定解区域的角点 $(0,0)$，$(l,0)$ 提了一些相容性条件。然而从物理学上的观点来看，这样的光滑性条件过于苛刻，也很不自然。我们举一个简单例子来阐明这个问题。考虑一根两端固定的弦，在中间的一点 $x_0$ 将弦线提起，使它离开水平位置，然后轻轻放下，弦就开始做横振动。这样的问题在物理学上是明确的，而在数学上如果我们拘泥于前面的解的定义，这样的问题在指定的函数类是无解的。因为此定解问题中初速度为零，初位移 $\varphi(x)$ 在点 $x_0$ 不可微（见图4.8），因此不满足古典解的存在性条件。实际上，由波动方程的性质，解在初值的奇性沿特征线向内部传播。更精确地说，当 $t>0$ 时，$\varphi'(x)$ 在点 $x_0$ 的不连续性将沿着特征线 $x-x_0=at$ 和 $x-x_0=-at$，并通过固定边界 $x=0,x=l$ 的反射，向定解区域内部传播开来。所以混合问题 \eqref{eq:wave-gen-mixed} 在 $Q$ 上不可能存在二次连续可微的解。

\begin{figure}[H]
\centering
\includegraphics[width=0.5\textwidth]{PDE图7.png}
\caption{初位移在点$x_0$不可微}
\label{figure:wave_initial_sing}
\end{figure}

为了使得这样的定解问题在数学上存在解，我们必须在一个更大的函数空间来求解，从而需要扩大解的函数类。因此我们必须放弃对解的二次连续可微的要求。我们通常将原来二次连续可微的解称为\textbf{古典解}。在扩大的函数类中的解称为\textbf{广义解}。为了使得我们在扩大的函数类找到的解符合实际问题的要求，我们在扩大解的函数类时必须遵守一定的原则。这些原则就是：

\begin{enumerate}[(1)]
\item 古典解必是广义解；
\item 广义解是唯一的，且在一定意义下连续依赖于定解条件。
\end{enumerate}

定义广义解的方式有很多种，但最自然的方式是回到原来的物理模型，按照守恒律或变分原理来定义。这是因为我们考虑的方程都是在假设解是二次连续可微的条件下，由守恒律的积分等式导出的。既然要降低解的光滑性要求，那么广义解满足方程只能在较弱的意义下理解。因此我们希望广义解在适当意义下满足守恒律的积分等式。

如果 $u(x,t)\in C^2(\overline{Q}_T)$ 是混合问题 \eqref{eq:wave-gen-mixed} 的古典解，对于任意 $\zeta(x,t)\in C^2(\overline{Q}_T)$，将混合问题 \eqref{eq:wave-gen-mixed} 的方程两边同乘 $\zeta(x,t)$，并在 $Q_T$ 上积分，则
\begin{align*}
\iint_{Q_T}(u_{tt}-a^2u_{xx})\zeta\,\mathrm{d}x\mathrm{d}t=0.
\end{align*}
对上述分部积分，利用边值条件将 $u$ 的偏导数转移到 $\zeta$ 上，得到
\begin{align*}
\iint_{Q_T}u(\zeta_{tt}-a^2\zeta_{xx})\,\mathrm{d}x\mathrm{d}t+\int_0^l(u_t\zeta-u\zeta_t)|_0^T\,\mathrm{d}x-a^2\int_0^T(u_x\zeta)|_0^l\,\mathrm{d}t=0.
\end{align*}
由 $u(x,t)$ 满足的初值条件，并假设 $\zeta(x,t)$ 满足
\begin{align*}
\zeta(x,T)=\zeta_t(x,T)=\zeta(0,t)=\zeta(l,t)=0,
\end{align*}
则上述简化为
\begin{align}
\iint_{Q_T}u(\zeta_{tt}-a^2\zeta_{xx})\,\mathrm{d}x\mathrm{d}t+\int_0^l\varphi\zeta_t(x,0)\,\mathrm{d}x-\int_0^l\psi\zeta(x,0)\,\mathrm{d}x=0. \label{eq:weak-form}
\end{align}
实际上，当 $u(x,t)\in C(\overline{Q}_T)$ 时，积分等式 \eqref{eq:weak-form} 有意义。这提示我们可以利用此积分等式来降低解的光滑性要求。

记
\begin{align*}
\mathcal{D}=\{\zeta(x,t)\in C^2(\overline{Q}_T)\mid \zeta(x,T)=\zeta_t(x,T)=\zeta(0,t)=\zeta(l,t)=0\}.
\end{align*}
$\mathcal{D}$ 称为\textbf{试验函数类}，$\zeta(x,t)\in\mathcal{D}$ 称为\textbf{试验函数}。

\begin{definition}
如果对于任意试验函数 $\zeta(x,t)\in\mathcal{D}$，积分等式 \eqref{eq:weak-form} 成立，则函数 $u(x,t)\in C(\overline{Q}_T)$ 称为\textbf{混合问题 \eqref{eq:wave-gen-mixed} 的广义解}。
\end{definition}

从广义解的引入过程中，我们知道古典解一定是广义解，即上述原则 (1) 成立。下面我们证明原则 (2) 成立。

\begin{theorem}
混合问题 \eqref{eq:wave-gen-mixed} 的广义解一定唯一。
\end{theorem}

\begin{proof}
设混合问题 \eqref{eq:wave-gen-mixed} 有两个解 $u_1,u_2$，并记 $w=u_1-u_2$。由于 $u_1,u_2$ 满足积分等式 \eqref{eq:weak-form}，则对于任意的 $\zeta\in\mathcal{D}$，函数 $w$ 满足
\begin{align}
\iint_{Q_T}w(\zeta_{tt}-a^2\zeta_{xx})\,\mathrm{d}x\mathrm{d}t=0. \label{eq:w-weak}
\end{align}
设 $g(x,t)\in C_0^\infty(Q_T)$，考虑定解问题
\begin{align}
\begin{cases}
\zeta_{tt}-a^2\zeta_{xx}=g(x,t), & (x,t)\in(0,l)\times(0,T), \\
\zeta(x,T)=0, & x\in[0,l], \\
\zeta_t(x,T)=0, & x\in[0,l], \\
\zeta(0,t)=\zeta(l,t)=0, & t\in[0,T].
\end{cases} \label{eq:dual-problem}
\end{align}
作变量替换 $\tau=T-t$ 和函数变换 $\overline{\zeta}(x,\tau)=\zeta(x,T-\tau)$，则 $\overline{\zeta}$ 在 $Q_T$ 上满足混合问题
\begin{align*}
\begin{cases}
\overline{\zeta}_{\tau\tau}-a^2\overline{\zeta}_{xx}=\overline{g}(x,\tau), & (x,\tau)\in(0,l)\times(0,T), \\
\overline{\zeta}(x,0)=0, & x\in[0,l], \\
\overline{\zeta}_\tau(x,0)=0, & x\in[0,l], \\
\overline{\zeta}(0,\tau)=\overline{\zeta}(l,\tau)=0, & \tau\in[0,T],
\end{cases}
\end{align*}
其中 $\overline{g}(x,\tau)=g(x,T-\tau)$。由于 $\overline{g}$ 是光滑函数且在 $\partial Q_T$ 的一个邻域上为 $0$，从而在角点满足相容性条件，由分离变量法知该混合问题存在解 $\overline{\zeta}\in C^2(\overline{Q}_T)$。显然函数 $\zeta(x,t)=\overline{\zeta}(x,T-t)\in\mathcal{D}$ 是定解问题 \eqref{eq:dual-problem} 的解。于是从 \eqref{eq:w-weak} 我们得到，对于任意 $g\in C_0^\infty(Q_T)$，成立
\begin{align}
\iint_{Q_T}wg\,\mathrm{d}x\mathrm{d}t=0. \label{eq:w-g}
\end{align}
下面证明在 $Q_T$ 上 $w\equiv 0$。用反证法。记 $z=(x,t)$。假设存在 $z_0=(x_0,t_0)\in Q_T$ 使得 $w(z_0)\neq 0$。不妨设 $w(z_0)>0$。由连续性，存在以 $z_0$ 为圆心、$r_0$ 为半径的圆盘 $B(z_0,r_0)\subset Q_T$，使得 $w(z)\geqslant \frac{1}{2}w(z_0)$。取 $g\in C_0^\infty(Q_T)$ 使得 $g\geqslant 0$，在 $B(z_0,\frac{r_0}{2})$ 上 $g\equiv 1$，在 $Q_T\setminus B(z_0,r_0)$ 上 $g\equiv 0$。于是
\begin{align*}
\iint_{Q_T}wg\,\mathrm{d}z
=\int_{B(z_0,r_0)}wg\,\mathrm{d}z
\geqslant \int_{B(z_0,\frac{r_0}{2})}wg\,\mathrm{d}z
\geqslant \frac{\pi r_0^2}{8}w(z_0)>0,
\end{align*}
这与 \eqref{eq:w-g} 矛盾，故 $w\equiv 0$。唯一性获证。
\end{proof}

\begin{theorem}
混合问题 \eqref{eq:wave-gen-mixed} 的广义解 $u(x,t)$ 满足
\begin{align}
\iint_{Q_T}u^2\,\mathrm{d}x\mathrm{d}t
\leqslant M\left(\int_0^l\varphi^2\,\mathrm{d}x+\int_0^l\psi^2\,\mathrm{d}x\right), \label{eq:energy-est-weak}
\end{align}
其中 $M$ 仅与 $a$ 和 $T$ 有关。
\end{theorem}

\begin{proof}
对于广义解 $u$，我们可以找到一列 $u_n\in C_0^\infty(Q_T)$ $(n=1,2,\cdots)$ 使得
\begin{align}
\iint_{Q_T}|u-u_n|^2\,\mathrm{d}x\mathrm{d}t\leqslant \frac{1}{n}. \label{eq:approx}
\end{align}
固定 $n$，考虑定解问题
\begin{align*}
\begin{cases}
\zeta_{tt}-a^2\zeta_{xx}=u_n, & (x,t)\in(0,l)\times(0,T), \\
\zeta(x,T)=0, & x\in[0,l], \\
\zeta_t(x,T)=0, & x\in[0,l], \\
\zeta(0,t)=\zeta(l,t)=0, & t\in[0,T].
\end{cases}
\end{align*}
由分离变量法知该问题存在解 $\zeta_n\in C^2(\overline{Q}_T)$。于是我们在积分等式 \eqref{eq:weak-form} 取试验函数 $\zeta_n$ 得到
\begin{align}
\iint_{Q_T}u\,u_n\,\mathrm{d}x\mathrm{d}t+\int_0^l\varphi(x)\zeta_{nt}(x,0)\,\mathrm{d}x-\int_0^l\psi(x)\zeta_n(x,0)\,\mathrm{d}x=0. \label{eq:weak-n}
\end{align}
由于 $\zeta_n$ 是古典解，由能量不等式可得
\begin{align*}
\int_0^l|\zeta_{nt}(x,0)|^2\,\mathrm{d}x &\leqslant C_1\iint_{Q_T}u_n^2\,\mathrm{d}x\mathrm{d}t,\\
\int_0^l|\zeta_n(x,0)|^2\,\mathrm{d}x &\leqslant C_2\iint_{Q_T}u_n^2\,\mathrm{d}x\mathrm{d}t.
\end{align*}
于是从 \eqref{eq:weak-n} 和不等式 $2ab\leqslant \varepsilon a^2+\frac{1}{\varepsilon}b^2$ 得
\begin{align*}
\iint_{Q_T}u u_n\,\mathrm{d}x\mathrm{d}t
&\leqslant \left(\int_0^l\varphi^2\,\mathrm{d}x\right)^{1/2}\left(\int_0^l|\zeta_{nt}|^2\,\mathrm{d}x\right)^{1/2}
+\left(\int_0^l\psi^2\,\mathrm{d}x\right)^{1/2}\left(\int_0^l|\zeta_n|^2\,\mathrm{d}x\right)^{1/2} \\
&\leqslant \frac{1}{2}\iint_{Q_T}u_n^2\,\mathrm{d}x\mathrm{d}t+C_3\int_0^l[\varphi^2+\psi^2]\,\mathrm{d}x.
\end{align*}
令 $n\to\infty$，利用 \eqref{eq:approx} 得
\begin{align*}
\iint_{Q_T}u^2\,\mathrm{d}x\mathrm{d}t
\leqslant \frac{1}{2}\iint_{Q_T}u^2\,\mathrm{d}x\mathrm{d}t+C_3\int_0^l[\varphi^2+\psi^2]\,\mathrm{d}x,
\end{align*}
移项即得 \eqref{eq:energy-est-weak}。
\end{proof}

\begin{corollary}
若 $u_i(x,t)$ 是相应于初值为 $\varphi_i(x)$，$\psi_i(x)$ $(i=1,2)$ 的混合问题 \eqref{eq:wave-gen-mixed} 的广义解，则
\begin{align}
\iint_{Q_T}|u_1-u_2|^2\,\mathrm{d}x\mathrm{d}t
\leqslant M\left(\int_0^l|\varphi_1-\varphi_2|^2\,\mathrm{d}x+\int_0^l|\psi_1-\psi_2|^2\,\mathrm{d}x\right), \label{eq:stab-weak}
\end{align}
其中 $M$ 仅与 $a$ 和 $T$ 有关。
\end{corollary}
\begin{remark}
这说明广义解在 $L^2$ 模意义下对于初值连续依赖。
\end{remark}

通过引入广义解的概念，我们可以大大降低对于初值的光滑性以及在角点的相容性条件的要求，使得一些过去认为无解的问题，在广义解的意义下也存在解。

\begin{theorem}
若 $\varphi(x)\in C([0,l])$，$\varphi(0)=\varphi(l)=0$，$\varphi'(x)$ 和 $\psi(x)$ 在 $[0,l]$ 上分别除去有限个第一类间断点外连续，则定解问题 \eqref{eq:wave-gen-mixed} 存在唯一的广义解，且仍可表示成级数形式
\begin{align}
u(x,t)=\sum_{n=1}^{\infty}\left(A_n\cos\frac{na\pi}{l}t+B_n\sin\frac{na\pi}{l}t\right)\sin\frac{n\pi}{l}x, \label{eq:gen-sol-series}
\end{align}
其中
\begin{align*}
A_n=\frac{2}{l}\int_0^l\varphi(x)\sin\frac{n\pi}{l}x\,\mathrm{d}x,\qquad
B_n=\frac{2}{na\pi}\int_0^l\psi(x)\sin\frac{n\pi}{l}x\,\mathrm{d}x.
\end{align*}
\end{theorem}

\begin{proof}
只需证明存在性。根据关于 $\varphi,\psi$ 的假设，它们可以按特征函数系 $\{\sin\frac{n\pi}{l}x\}$ 展开为
\begin{align*}
\varphi(x)=\sum_{n=1}^{\infty}\varphi_n\sin\frac{n\pi}{l}x,\qquad
\psi(x)=\sum_{n=1}^{\infty}\psi_n\sin\frac{n\pi}{l}x,
\end{align*}
其中 $\varphi_n,\psi_n$ 是 Fourier 系数。由 $\varphi$ 满足的条件，可得 $\varphi_n=\frac{l}{n\pi}\varphi'_n$，其中 $\varphi'_n$ 是 $\varphi'(x)$ 按 $\{\cos\frac{n\pi}{l}x\}$ 展开的 Fourier 系数。记
\begin{align*}
S_N(x)=\sum_{n=0}^{N}\varphi'_n\cos\frac{n\pi}{l}x,\qquad
T_N(x)=\sum_{n=1}^{N}\psi_n\sin\frac{n\pi}{l}x.
\end{align*}
由特征函数系的完备性，有
\begin{align}
\lim_{N\to\infty}\int_0^l|\varphi'(x)-S_N(x)|^2\,\mathrm{d}x=0, \label{eq:phi-approx}\\
\lim_{N\to\infty}\int_0^l|\psi(x)-T_N(x)|^2\,\mathrm{d}x=0. \label{eq:psi-approx}
\end{align}
由 Bessel 不等式可得
\begin{align}
\sum_{n=1}^{\infty}(\varphi'_n)^2\leqslant \int_0^l|\varphi'(x)|^2\,\mathrm{d}x,\qquad
\sum_{n=1}^{\infty}\psi_n^2\leqslant \int_0^l|\psi(x)|^2\,\mathrm{d}x. \label{eq:bessel}
\end{align}
考虑混合问题
\begin{align*}
\begin{cases}
u_{Ntt}-a^2u_{Nxx}=0, & (x,t)\in Q_T, \\
u_N(x,0)=S_N(x), & x\in[0,l], \\
u_{Nt}(x,0)=T_N(x), & x\in[0,l], \\
u_N(0,t)=u_N(l,t)=0, & t\in[0,T],
\end{cases}
\end{align*}
由古典解的存在性定理，其解可表为
\begin{align*}
u_N(x,t)=\sum_{n=1}^{N}\left(A_n\cos\frac{na\pi}{l}t+B_n\sin\frac{na\pi}{l}t\right)\sin\frac{n\pi}{l}x,
\end{align*}
其中 $A_n=\frac{l}{n\pi}\varphi'_n$，$B_n=\frac{l}{n\pi}\psi_n$。显然 $u_N$ 是级数 \eqref{eq:gen-sol-series} 的前 $N$ 项部分和。在 $u_N$ 的方程两边乘以 $\zeta\in\mathcal{D}$，分部积分得
\begin{align}
\iint_{Q_T}u_N(\zeta_{tt}-a^2\zeta_{xx})\,\mathrm{d}x\mathrm{d}t+\int_0^l S_N(x)\zeta_t(x,0)\,\mathrm{d}x-\int_0^l T_N(x)\zeta(x,0)\,\mathrm{d}x=0. \label{eq:weak-N}
\end{align}
令 $u$ 为级数 \eqref{eq:gen-sol-series} 的和函数。先证 $u_N$ 一致收敛到 $u$。记
\begin{align*}
r_N=u-u_N=\sum_{n=N+1}^{\infty}\left(A_n\cos\frac{na\pi}{l}t+B_n\sin\frac{na\pi}{l}t\right)\sin\frac{n\pi}{l}x.
\end{align*}
则
\begin{align*}
|r_N|&\leqslant \sum_{n=N+1}^{\infty}(|A_n|+|B_n|)
=\sum_{n=N+1}^{\infty}\frac{l}{n\pi}\left(|\varphi'_n|+\frac{1}{a}|\psi_n|\right) \\
&\leqslant \sum_{n=N+1}^{\infty}\left(\frac{l}{n\pi}\right)^2+\frac{1}{2}\sum_{n=N+1}^{\infty}\left(|\varphi'_n|^2+\frac{1}{a^2}|\psi_n|^2\right),
\end{align*}
由 \eqref{eq:bessel} 知当 $N\to\infty$ 时右端趋于 $0$，故 $u_N$ 一致收敛到 $u$。由于 $u_N\in C(\overline{Q}_T)$，故 $u\in C(\overline{Q}_T)$。在 \eqref{eq:weak-N} 中令 $N\to\infty$，利用 \eqref{eq:phi-approx} 和 \eqref{eq:psi-approx} 可得
\begin{align*}
\iint_{Q_T}u(\zeta_{tt}-a^2\zeta_{xx})\,\mathrm{d}x\mathrm{d}t+\int_0^l\varphi(x)\zeta_t(x,0)\,\mathrm{d}x-\int_0^l\psi(x)\zeta(x,0)\,\mathrm{d}x=0,
\end{align*}
即 $u$ 满足广义解定义。存在性得证。唯一性已由定理给出。
\end{proof}

\begin{remark}
定理和推论实际上说明定解问题 \eqref{eq:wave-gen-mixed} 在广义解的框架下仍然是适定的。
\end{remark}


\end{document}